\documentclass[12pt]{article}
\usepackage[margin=0.8in]{geometry}
\usepackage{amsmath}
\usepackage{xcolor}
\usepackage{tcolorbox}
\usepackage{enumitem}

\title{\textbf{Portfolio Choice: Understanding the FOC and Corners}}
\author{}
\date{}

\begin{document}
\maketitle

\section*{The Problem Setup}

You have wealth $w_0$ to invest. You can:
\begin{itemize}
\item Invest amount $x$ in a risky asset
\item Keep the rest $(w_0 - x)$ safe (earning nothing)
\end{itemize}

The risky asset:
\begin{itemize}
\item Pays off with probability $p$ (you get your money back plus a return)
\item Fails with probability $(1-p)$ (you lose what you invested)
\end{itemize}

Your utility is CRRA: $u(w) = \frac{w^{1-\sigma}}{1-\sigma}$ where $\sigma$ is your risk aversion.

\textbf{Goal:} Choose $x$ (how much to invest in risky asset) to maximize expected utility.

\section*{What the FOC Tells You}

When you take the first-order condition (derivative = 0) and solve, you get some formula for $x^*$.

\begin{tcolorbox}[colback=blue!5!white,colframe=blue!75!black]
\textbf{The Problem:}

The FOC assumes you're at an \textbf{interior solution} where $0 < x < w_0$.

But sometimes the "optimal" $x$ from the FOC is:
\begin{itemize}
\item Negative (invest less than zero = borrow to invest)
\item Bigger than $w_0$ (invest more than you have)
\end{itemize}

These violate the constraint that $0 \leq x \leq w_0$!
\end{tcolorbox}

So after finding $x^*$ from the FOC, you need to \textbf{check corners}:
\begin{enumerate}
\item Is $x^* \geq 0$? If not, set $x = 0$
\item Is $x^* \leq w_0$? If not, set $x = w_0$
\end{enumerate}

\newpage

\section*{Breaking Down the Passage}

\subsection*{Part 1: "FOC is only valid if $\sigma > 0$"}

\begin{tcolorbox}[colback=yellow!10!white,colframe=orange!75!black]
\textbf{If $\sigma = 0$ (risk-neutral):}

You only care about expected value, not risk. So your decision is simple:

\begin{itemize}
\item If the risky asset has positive expected return: invest ALL your wealth ($x = w_0$)
\item If the risky asset has negative expected return: invest NOTHING ($x = 0$)
\end{itemize}

There's no interior solution! You go all-in or stay out completely.

The FOC doesn't apply because the optimal choice is always at a corner.
\end{tcolorbox}

\textbf{Why?} Risk-neutral people don't care about diversification. They just chase the highest expected return.

\subsection*{Part 2: "FOC is valid provided $p \geq \frac{1}{2}$"}

\begin{tcolorbox}[colback=green!5!white,colframe=green!75!black]
\textbf{The FOC gives you a formula for $x^*$. But we need $x^* \geq 0$ for it to make sense.}

When you work out the math (the actual FOC), you find that:
\[
x^* \geq 0 \quad \Leftrightarrow \quad p \geq \frac{1}{2}
\]

In other words:
\begin{itemize}
\item If $p \geq \frac{1}{2}$: The FOC gives a valid interior solution with $0 < x^* < w_0$
\item If $p < \frac{1}{2}$: The FOC would give $x^* < 0$ (negative investment), which is impossible!
\end{itemize}
\end{tcolorbox}

\subsection*{Part 3: "If $p > \frac{1}{2}$, invest amount strictly positive and strictly below $w_0$"}

This means:
\begin{itemize}
\item The risky asset is "good enough" (succeeds more than half the time)
\item You should invest some of your wealth: $x^* > 0$
\item But not all of it: $x^* < w_0$ (because you're risk-averse and want to diversify)
\end{itemize}

\textbf{Example:} If $p = 0.6$ (60\% chance of success), $\sigma = 2$, $w_0 = 100$:

The FOC might give something like $x^* = 40$. You invest \$40 in the risky asset and keep \$60 safe.

\newpage

\subsection*{Part 4: "If $p < \frac{1}{2}$, best option is to invest zero"}

\begin{tcolorbox}[colback=red!5!white,colframe=red!75!black]
\textbf{If the risky asset fails more than half the time:}

The FOC would tell you to invest a negative amount (which means borrow money to short the asset). But you can't do that!

Given the constraint $x \geq 0$, your best option is: \textbf{don't invest at all}.

Set $x = 0$ (corner solution).
\end{tcolorbox}

\textbf{Intuition:} If $p < 0.5$, the risky asset is a bad bet on average. Since you can't short it, you just avoid it completely.

\subsection*{Part 5: "Optimal investment increases in $p$, decreases in $\sigma$"}

\textbf{Increases in $p$:}
\begin{itemize}
\item Higher $p$ = risky asset is more likely to pay off
\item So you invest more: $x^*$ goes up
\end{itemize}

\textbf{Example:}
\begin{itemize}
\item If $p = 0.6$: maybe you invest \$40
\item If $p = 0.7$: you invest \$60 (more confident it will pay off)
\end{itemize}

\vspace{0.3cm}

\textbf{Decreases in $\sigma$:}
\begin{itemize}
\item Higher $\sigma$ = more risk-averse
\item So you invest less in the risky asset: $x^*$ goes down
\end{itemize}

\textbf{Example:}
\begin{itemize}
\item If $\sigma = 1$ (moderately risk-averse): invest \$50
\item If $\sigma = 3$ (very risk-averse): invest \$20 (prefer safety)
\end{itemize}

\newpage

\section*{The Big Picture: Checking Corners}

Here's the full decision tree:

\begin{center}
\fbox{\begin{minipage}{0.9\textwidth}
\textbf{Step 1: Check if risk-averse}

Is $\sigma > 0$?

\vspace{0.2cm}

\textbf{No ($\sigma = 0$):} You're risk-neutral
\begin{itemize}
\item Calculate expected return of risky asset
\item If positive: $x^* = w_0$ (invest everything)
\item If negative: $x^* = 0$ (invest nothing)
\end{itemize}

\textbf{Yes ($\sigma > 0$):} Go to Step 2

\vspace{0.5cm}

\textbf{Step 2: Check if risky asset is worth it}

Is $p \geq \frac{1}{2}$?

\vspace{0.2cm}

\textbf{No ($p < \frac{1}{2}$):} Risky asset is a bad bet
\begin{itemize}
\item Corner solution: $x^* = 0$ (invest nothing)
\end{itemize}

\textbf{Yes ($p \geq \frac{1}{2}$):} Go to Step 3

\vspace{0.5cm}

\textbf{Step 3: Use the FOC}

Solve the FOC to get $x^*$

\vspace{0.2cm}

This gives an interior solution: $0 < x^* < w_0$

\vspace{0.2cm}

Check comparative statics:
\begin{itemize}
\item Higher $p$ $\Rightarrow$ higher $x^*$ (invest more)
\item Higher $\sigma$ $\Rightarrow$ lower $x^*$ (invest less)
\end{itemize}

\end{minipage}}
\end{center}

\newpage

\section*{Concrete Example}

Let's say the risky asset:
\begin{itemize}
\item Returns \$2 for every \$1 invested with probability $p$
\item Returns \$0 with probability $(1-p)$
\end{itemize}

You have $w_0 = \$100$ and utility $u(w) = \frac{w^{1-\sigma}}{1-\sigma}$.

\subsection*{Case 1: $p = 0.4$ (risky asset fails 60\% of the time)}

Since $p < 0.5$, the risky asset is a bad bet.

\textbf{Decision:} Don't invest at all. $x^* = 0$.

Keep all \$100 safe.

\subsection*{Case 2: $p = 0.6$, $\sigma = 2$ (60\% success, moderately risk-averse)}

Since $p > 0.5$ and $\sigma > 0$, use the FOC.

The FOC would give some formula like:
\[
x^* = f(p, \sigma, w_0)
\]

Let's say it gives $x^* = \$40$.

\textbf{Decision:} Invest \$40 in risky asset, keep \$60 safe.

\textbf{Expected outcome:}
\begin{itemize}
\item 60\% chance: risky asset pays \$80, plus \$60 safe = \$140 total
\item 40\% chance: risky asset pays \$0, plus \$60 safe = \$60 total
\end{itemize}

\subsection*{Case 3: Same but $\sigma = 4$ (very risk-averse)}

FOC might now give $x^* = \$20$.

\textbf{Decision:} Invest only \$20 in risky asset (less than before because you're more scared of risk).

\section*{Summary}

\begin{tcolorbox}[colback=green!5!white,colframe=green!75!black]
\textbf{Key Takeaways:}

\begin{enumerate}[leftmargin=*]
\item \textbf{FOC only works if $\sigma > 0$ (risk-averse)}
\begin{itemize}
\item If risk-neutral ($\sigma = 0$): go all-in or stay out
\end{itemize}

\item \textbf{FOC gives valid $x^*$ only if $p \geq \frac{1}{2}$}
\begin{itemize}
\item If $p < \frac{1}{2}$: corner solution, set $x = 0$
\end{itemize}

\item \textbf{Interior solution when $p > \frac{1}{2}$ and $\sigma > 0$}
\begin{itemize}
\item Invest some but not all: $0 < x^* < w_0$
\end{itemize}

\item \textbf{Comparative statics:}
\begin{itemize}
\item $x^*$ increases in $p$ (better odds $\Rightarrow$ invest more)
\item $x^*$ decreases in $\sigma$ (more risk-averse $\Rightarrow$ invest less)
\end{itemize}

\item \textbf{Always check corners!}
\begin{itemize}
\item Is $x^* \geq 0$? If not, set $x = 0$
\item Is $x^* \leq w_0$? If not, set $x = w_0$
\end{itemize}
\end{enumerate}
\end{tcolorbox}

\end{document}
